% Options for packages loaded elsewhere
\PassOptionsToPackage{unicode}{hyperref}
\PassOptionsToPackage{hyphens}{url}
\documentclass[
]{article}
\usepackage{xcolor}
\usepackage{amsmath,amssymb}
\setcounter{secnumdepth}{-\maxdimen} % remove section numbering
\usepackage{iftex}
\ifPDFTeX
  \usepackage[T1]{fontenc}
  \usepackage[utf8]{inputenc}
  \usepackage{textcomp} % provide euro and other symbols
\else % if luatex or xetex
  \usepackage{unicode-math} % this also loads fontspec
  \defaultfontfeatures{Scale=MatchLowercase}
  \defaultfontfeatures[\rmfamily]{Ligatures=TeX,Scale=1}
\fi
\usepackage{lmodern}
\ifPDFTeX\else
  % xetex/luatex font selection
\fi
% Use upquote if available, for straight quotes in verbatim environments
\IfFileExists{upquote.sty}{\usepackage{upquote}}{}
\IfFileExists{microtype.sty}{% use microtype if available
  \usepackage[]{microtype}
  \UseMicrotypeSet[protrusion]{basicmath} % disable protrusion for tt fonts
}{}
\makeatletter
\@ifundefined{KOMAClassName}{% if non-KOMA class
  \IfFileExists{parskip.sty}{%
    \usepackage{parskip}
  }{% else
    \setlength{\parindent}{0pt}
    \setlength{\parskip}{6pt plus 2pt minus 1pt}}
}{% if KOMA class
  \KOMAoptions{parskip=half}}
\makeatother
\usepackage{color}
\usepackage{fancyvrb}
\newcommand{\VerbBar}{|}
\newcommand{\VERB}{\Verb[commandchars=\\\{\}]}
\DefineVerbatimEnvironment{Highlighting}{Verbatim}{commandchars=\\\{\}}
% Add ',fontsize=\small' for more characters per line
\newenvironment{Shaded}{}{}
\newcommand{\AlertTok}[1]{\textcolor[rgb]{1.00,0.00,0.00}{\textbf{#1}}}
\newcommand{\AnnotationTok}[1]{\textcolor[rgb]{0.38,0.63,0.69}{\textbf{\textit{#1}}}}
\newcommand{\AttributeTok}[1]{\textcolor[rgb]{0.49,0.56,0.16}{#1}}
\newcommand{\BaseNTok}[1]{\textcolor[rgb]{0.25,0.63,0.44}{#1}}
\newcommand{\BuiltInTok}[1]{\textcolor[rgb]{0.00,0.50,0.00}{#1}}
\newcommand{\CharTok}[1]{\textcolor[rgb]{0.25,0.44,0.63}{#1}}
\newcommand{\CommentTok}[1]{\textcolor[rgb]{0.38,0.63,0.69}{\textit{#1}}}
\newcommand{\CommentVarTok}[1]{\textcolor[rgb]{0.38,0.63,0.69}{\textbf{\textit{#1}}}}
\newcommand{\ConstantTok}[1]{\textcolor[rgb]{0.53,0.00,0.00}{#1}}
\newcommand{\ControlFlowTok}[1]{\textcolor[rgb]{0.00,0.44,0.13}{\textbf{#1}}}
\newcommand{\DataTypeTok}[1]{\textcolor[rgb]{0.56,0.13,0.00}{#1}}
\newcommand{\DecValTok}[1]{\textcolor[rgb]{0.25,0.63,0.44}{#1}}
\newcommand{\DocumentationTok}[1]{\textcolor[rgb]{0.73,0.13,0.13}{\textit{#1}}}
\newcommand{\ErrorTok}[1]{\textcolor[rgb]{1.00,0.00,0.00}{\textbf{#1}}}
\newcommand{\ExtensionTok}[1]{#1}
\newcommand{\FloatTok}[1]{\textcolor[rgb]{0.25,0.63,0.44}{#1}}
\newcommand{\FunctionTok}[1]{\textcolor[rgb]{0.02,0.16,0.49}{#1}}
\newcommand{\ImportTok}[1]{\textcolor[rgb]{0.00,0.50,0.00}{\textbf{#1}}}
\newcommand{\InformationTok}[1]{\textcolor[rgb]{0.38,0.63,0.69}{\textbf{\textit{#1}}}}
\newcommand{\KeywordTok}[1]{\textcolor[rgb]{0.00,0.44,0.13}{\textbf{#1}}}
\newcommand{\NormalTok}[1]{#1}
\newcommand{\OperatorTok}[1]{\textcolor[rgb]{0.40,0.40,0.40}{#1}}
\newcommand{\OtherTok}[1]{\textcolor[rgb]{0.00,0.44,0.13}{#1}}
\newcommand{\PreprocessorTok}[1]{\textcolor[rgb]{0.74,0.48,0.00}{#1}}
\newcommand{\RegionMarkerTok}[1]{#1}
\newcommand{\SpecialCharTok}[1]{\textcolor[rgb]{0.25,0.44,0.63}{#1}}
\newcommand{\SpecialStringTok}[1]{\textcolor[rgb]{0.73,0.40,0.53}{#1}}
\newcommand{\StringTok}[1]{\textcolor[rgb]{0.25,0.44,0.63}{#1}}
\newcommand{\VariableTok}[1]{\textcolor[rgb]{0.10,0.09,0.49}{#1}}
\newcommand{\VerbatimStringTok}[1]{\textcolor[rgb]{0.25,0.44,0.63}{#1}}
\newcommand{\WarningTok}[1]{\textcolor[rgb]{0.38,0.63,0.69}{\textbf{\textit{#1}}}}
\setlength{\emergencystretch}{3em} % prevent overfull lines
\providecommand{\tightlist}{%
  \setlength{\itemsep}{0pt}\setlength{\parskip}{0pt}}
\usepackage{bookmark}
\IfFileExists{xurl.sty}{\usepackage{xurl}}{} % add URL line breaks if available
\urlstyle{same}
\hypersetup{
  pdftitle={EV Adoption Model Literature Review},
  hidelinks,
  pdfcreator={LaTeX via pandoc}}

\title{EV Adoption Model Literature Review}
\author{}
\date{}

\begin{document}
\maketitle

\subsection{High-Level Findings}\label{high-level-findings}

EV adoption is usually not modeled as a single consumer-choice event.
The stronger models combine several mechanisms: household heterogeneity,
price and total-cost barriers, charging access, social visibility,
policy incentives, technology learning, and vehicle replacement timing.
For an affordance model, the most transferable idea is to treat EV
adoption as the result of both objective action opportunities and
perceived feasibility.

The closest analogue to the current affordance model is the agent-based
EV adoption literature. Studies such as Eppstein et al.~on PHEV market
penetration, McCoy and Lyons on Irish microsimulation, Brown's
mixed-logit ABM, and Kangur et al.'s psychologically grounded ABM all
represent heterogeneous agents who update or act under social and
contextual influence. Useful sources include
\href{https://doi.org/10.1016/j.enpol.2011.04.007}{Eppstein et
al.~2011}, \href{https://doi.org/10.1016/j.erss.2014.07.008}{McCoy and
Lyons 2014}, \href{https://doi.org/10.18564/jasss.2127}{Brown 2013}, and
\href{https://doi.org/10.1016/j.jenvp.2017.01.002}{Kangur et al.~2017}.

Aggregate diffusion and system dynamics models are useful for feedback
structure, not micro-behavioral detail. They emphasize reinforcing loops
among adoption, infrastructure, cost reductions, legitimacy, and policy
support. Representative sources include
\href{https://doi.org/10.1068/b33022t}{Struben and Sterman 2008},
\href{https://doi.org/10.1016/j.tranpol.2011.12.006}{Shepherd et
al.~2012},
\href{https://doi.org/10.1016/j.techfore.2015.11.028}{Pasaoglu et
al.~2016}, and the review by
\href{https://doi.org/10.1016/j.rser.2018.03.055}{Gnann et al.~2018}.

Discrete choice and hybrid choice models provide the strongest empirical
foundation for purchase probabilities. They estimate willingness to pay
for price, range, charging time, fuel/electricity cost, infrastructure,
incentives, and latent attitudes. Representative work includes
Brownstone, Bunch, and Train on joint RP/SP mixed logit, Hidrue et
al.~on willingness to pay, Glerum et al. on hybrid choice, Helveston et
al.~on U.S. and China preferences, and Axsen et al.~on lifestyle
heterogeneity. Relevant links include
\href{https://www.sciencedirect.com/science/article/pii/S0928765511000200}{Hidrue
et al.~2011}, \href{https://doi.org/10.1287/trsc.2013.0487}{Glerum et
al.~2014}, and
\href{https://www.cmu.edu/me/ddl/publications/2015-TRA-Helveston-etal-EVs-in-China-US.pdf}{Helveston
et al.~2015}.

Social network and peer-effect studies show that EV uptake is socially
interdependent. Peer adoption, local visibility, word of mouth, and
perceived normality shift the effective choice set. This maps directly
onto social learning in the affordance model. Useful sources include
\href{https://www.sciencedirect.com/science/article/pii/S2214629614000863}{McCoy
and Lyons 2014},
\href{https://research.rug.nl/en/publications/an-agent-based-model-for-diffusion-of-electric-vehicles}{Kangur
et al.~2017},
\href{https://www.sciencedirect.com/science/article/pii/S1361920916301031}{Liu
et al.~2017}, and
\href{https://dspace.library.uu.nl/handle/1874/375563}{Edelenbosch et
al.~2018}.

Charging infrastructure and spatial adoption models are central for an
affordance framing. Charging availability is not just an explanatory
variable; it is a local opportunity structure. Stronger studies
distinguish charger counts, accessibility, fast-charging corridors, home
charging, land-use context, and spatial spillovers. Representative
sources include
\href{https://doi.org/10.1016/j.enpol.2014.01.043}{Sierzchula et
al.~2014}, \href{https://doi.org/10.1086/689702}{Li et al.~2017},
\href{https://doi.org/10.1088/1748-9326/aad0f8}{Narassimhan and Johnson
2018}, \href{https://doi.org/10.1016/j.erss.2022.102663}{White et
al.~2022}, and \href{https://doi.org/10.1016/j.trd.2024.104400}{Qian et
al.~2024}.

Policy and total cost of ownership models are the best source for cost
mechanisms. They translate purchase subsidies, taxes, fuel prices,
electricity prices, charging costs, depreciation, maintenance, and
vehicle lifetime into adoption conditions. Useful sources include
\href{https://doi.org/10.1016/j.enpol.2016.03.050}{Langbroek et
al.~2016}, \href{https://doi.org/10.1016/j.enpol.2018.04.024}{Danielis
et al.~2018}, \href{https://doi.org/10.1016/j.eneco.2019.104493}{Munzel
et al.~2019}, and
\href{https://doi.org/10.1016/j.apenergy.2024.124838}{Noll et al.~2024}.

\subsection{Transferable Mechanisms For An Affordance-Based EV
ABM}\label{transferable-mechanisms-for-an-affordance-based-ev-abm}

The most useful model design is modular:

\begin{enumerate}
\def\labelenumi{\arabic{enumi}.}
\tightlist
\item
  Affordance layer: local charging access, home charging, public
  charging, fast-charging corridors, dealership/model availability,
  parking access, and financial affordability.
\item
  Personal-state layer: environmental orientation, innovation affinity,
  risk tolerance, range anxiety, price sensitivity, and EV familiarity.
\item
  Choice layer: probability of EV purchase conditional on replacement
  need, affordability, charging feasibility, and perceived usefulness.
\item
  Social layer: peer exposure, neighborhood visibility, word of mouth,
  and social norm updating.
\item
  Policy layer: purchase subsidies, taxes, fuel/electricity prices,
  charging infrastructure support, and information campaigns.
\item
  Feedback layer: adoption increases visibility, encourages charger
  rollout, reduces perceived risk, and may reduce costs through
  technology learning.
\end{enumerate}

For the existing Mesa affordance model, the most natural extension is
not to replace the behavior rule with a full vehicle-choice model
immediately. A safer first step is to add EV-specific affordance fields
and let them modify the probability of pro-EV behavior:

\begin{Shaded}
\begin{Highlighting}[]
\NormalTok{effective\_EV\_propensity =}
\NormalTok{    personal\_EV\_state}
\NormalTok{    + peer\_exposure\_effect}
\NormalTok{    + charging\_affordance\_effect}
\NormalTok{    {-} price\_or\_TCO\_barrier}
\NormalTok{    {-} range\_anxiety\_barrier}
\end{Highlighting}
\end{Shaded}

This preserves the affordance logic while allowing empirical mechanisms
from EV adoption models to enter as interpretable modules.

\subsection{Specific Design Guidance}\label{specific-design-guidance}

For price mechanisms, use TCO literature to represent cost as more than
upfront price. Include purchase price, subsidy, expected fuel savings,
electricity or charging price, annual mileage, maintenance, home charger
cost, and depreciation where data are available. Price should affect
feasibility and adoption probability, not simply subtract from
environmental attitude.

For infrastructure, avoid using only a global charger count. Treat
charging as a spatial affordance with home access, nearby public access,
destination access, fast-charging corridor access, reliability, and
congestion if possible.

For social learning, use local adoption exposure or network-neighbor
adoption to change perceived feasibility, risk, and normality. Peer
effects should be allowed to amplify infrastructure and policy effects.

For heterogeneity, represent households by at least income or budget,
dwelling type, annual mileage, car replacement timing, environmental
preference, and innovation affinity. These are repeatedly important
across choice, ABM, TCO, and infrastructure studies.

\subsection{Key Gaps In The
Literature}\label{key-gaps-in-the-literature}

Several gaps are relevant for a new research contribution:

\begin{enumerate}
\def\labelenumi{\arabic{enumi}.}
\tightlist
\item
  Many models still weakly connect micro-level social influence with
  spatial charging access.
\item
  Used EV markets, resale value, vehicle availability, and supply
  constraints are under-modeled.
\item
  Home charging, renter status, multifamily housing, and equity
  constraints are often too coarse.
\item
  Public charging is often represented as station count, not
  reliability, congestion, price, or practical accessibility.
\item
  Preference updating after trial, ownership experience, and peer
  information remains simplified.
\item
  Policy interactions are usually clearer in system dynamics and TCO
  models than in household-level ABMs.
\end{enumerate}

\end{document}
